\chapter{ASR --- Reconnaissance Automatique de la Parole}

\section{Identification}

\begin{center}
\begin{tabular}{ll}
\toprule
\textbf{Champ} & \textbf{Valeur} \\
\midrule
Modèle         & \texttt{openai/whisper-large-v3} \\
Mode           & Chunked (30\,s/chunk, batch = 8) \\
Statut scoring & À implémenter \\
Fichier        & \texttt{DataPipelines/WhisperModel.py} \\
\bottomrule
\end{tabular}
\end{center}

\section{Fonctionnalité}

Whisper transcrit les fichiers audio et vidéo en texte. Dans le contexte
judiciaire : écoutes téléphoniques, auditions, vidéos de surveillance.

\section{Modèles évalués}

\begin{center}
\begin{tabular}{llll}
\toprule
\textbf{Modèle} & \textbf{Version} & \textbf{Benchmarké} & \textbf{Retenu} \\
\midrule
whisper-large-v3   & large-v3      & Non  & \textbf{Oui} \\
whisper-medium     & medium        & Non  & Non \\
Faster-Whisper CT2 & CT2 large-v3  & Non  & Non \\
\bottomrule
\end{tabular}
\end{center}

\section{Fonction de scoring de confiance}

\texttt{WhisperModel.infer()} retourne uniquement \texttt{result["text"]} ---
aucun beam score ni métadonnée de confiance n'est exposé dans l'état actuel.

\textbf{Pour implémenter un scoring}, il faudrait modifier \texttt{infer()} afin
de récupérer les \texttt{output\_scores} de la pipeline HuggingFace.
Une approche cohérente avec \texttt{TranslationModel} serait d'utiliser
les log-probs best-beam par chunk (même architecture seq2seq, même beam search).
Cette piste reste à valider --- aucune formule n'est fixée à ce stade.

\section{Justification scientifique}

Architecture Whisper~\cite{radford2022whisper}.

\section{Validation empirique}

Aucun dataset d'évaluation annoté n'est disponible à ce jour. Le
\texttt{compute\_confidence()} est actuellement validé par inspection manuelle
sur des cas réels et des scénarios limites (\textit{smoke testing}).

\noindent\textbf{Métriques à mesurer lors de la constitution du dataset :}
\begin{itemize}
  \item Corrélation entre le score de confiance et la métrique de qualité
    appropriée (ici $1 - \mathrm{WER}$) ;
  \item Calibration (ECE) : un score de confiance de 0.8 correspond-il à
    environ 80\,\% de réussite sur le dataset ?
  \item Pouvoir discriminant (AUC-ROC) : le score sépare-t-il les transcriptions
    correctes des incorrectes ?
\end{itemize}

\noindent\textbf{Statut :} \textit{À implémenter.}
